\clearpage
\label{s:app:start}
\begin{center}{\fontsize{16}{18}\selectfont\bfseries APPENDIX}\end{center}\vspace{2pt}
{\fontsize{14}{16}\selectfont\bfseries A.\ \ Full Forms of Abbreviations}\par\vspace{4pt}
\begingroup
\fontsize{11}{12.5}\selectfont\renewcommand{\baselinestretch}{1}\selectfont\setlength{\parskip}{0pt}
\renewcommand{\arraystretch}{1.05}
\setlength{\LTleft}{0pt}\setlength{\LTright}{0pt}
\begin{longtable}{|p{110pt}|p{343pt}|}
\hline \textbf{Abbreviation} & \textbf{Full form}\\
\endhead
\hline ACMG/AMP & American College of Medical Genetics and Genomics / Association for Molecular Pathology\\
\hline AI & Artificial Intelligence\\
\hline AI \& DS & Artificial Intelligence and Data Science\\
\hline API & Application Programming Interface\\
\hline AUROC & Area Under the Receiver Operating Characteristic curve\\
\hline bp / kb & base pair / kilobase (1,000 base pairs)\\
\hline BRCA1 & Breast Cancer gene 1\\
\hline CADD & Combined Annotation Dependent Depletion\\
\hline ClinGen & Clinical Genome Resource\\
\hline ClinVar & Clinical Variation database of the NCBI\\
\hline CSS & Cascading Style Sheets\\
\hline CUDA & Compute Unified Device Architecture\\
\hline DBIT & Don Bosco Institute of Technology\\
\hline DNA & Deoxyribonucleic Acid\\
\hline GB / GHz & Gigabyte / Gigahertz\\
\hline GPN & Genomic Pre-trained Network\\
\hline GPU & Graphics Processing Unit\\
\hline GRCh37 / GRCh38 & Genome Reference Consortium Human genome builds 37 and 38\\
\hline HGVS & Human Genome Variation Society\\
\hline HS256 & HMAC (Hash-based Message Authentication Code) using SHA-256\\
\hline HTTP / HTTPS & Hypertext Transfer Protocol / Hypertext Transfer Protocol Secure\\
\hline JSON & JavaScript Object Notation\\
\hline JWT & JSON Web Token\\
\hline LOF & Loss of Function\\
\hline MSA & Multiple Sequence Alignment\\
\hline NCBI & National Center for Biotechnology Information\\
\hline NLM & National Library of Medicine\\
\hline PolyPhen-2 & Polymorphism Phenotyping, version 2\\
\hline RAM & Random Access Memory\\
\hline REST & Representational State Transfer\\
\hline ROC & Receiver Operating Characteristic\\
\hline SGE & Saturation Genome Editing\\
\hline SIFT & Sorting Intolerant From Tolerant\\
\hline SNV & Single Nucleotide Variant\\
\hline SSD & Solid State Drive\\
\hline UCSC & University of California, Santa Cruz\\
\hline UI & User Interface\\
\hline URL & Uniform Resource Locator\\
\hline VCF & Variant Call Format\\
\hline VTU & Visvesvaraya Technological University\\
\hline
\end{longtable}
\endgroup
\clearpage
{\fontsize{14}{16}\selectfont\bfseries B.\ \ Formulas and Notation}\par\vspace{4pt}
\begingroup
\fontsize{12}{14}\selectfont\renewcommand{\baselinestretch}{1.15}\selectfont\setlength{\parskip}{4pt}
\renewcommand{\theequation}{A.\arabic{equation}}\setcounter{equation}{0}
The formulas below are the ones used by the service and by the analysis in this report. Here $p$ is the 1-based position of the variant, $W=8192$ is the nominal window size, $w^{ref}$ and $w^{var}$ are the reference and the mutated window, and $S(\cdot)$ is the Evo~2-7B log-likelihood score of a sequence.

\textbf{Genome window} (half-open, 0-based) and the offset of the variant in it:
\begin{equation}\mathrm{start}=\max\!\left(0,\;p-1-\tfrac{W}{2}\right),\qquad \mathrm{end}=p+\tfrac{W}{2},\qquad \mathrm{offset}=p-1-\mathrm{start}\end{equation}
This gives $W+1=8193$ bases with the variant at offset 4,096.

\textbf{Score difference:}
\begin{equation}\Delta=S(w^{var})-S(w^{ref})\end{equation}

\textbf{Sensitivity, specificity and false-positive rate} (loss of function is the positive class):
\begin{equation}\mathrm{TPR}=\frac{TP}{TP+FN},\qquad \mathrm{TNR}=\frac{TN}{TN+FP},\qquad \mathrm{FPR}=\frac{FP}{FP+TN}=1-\mathrm{TNR}\end{equation}

\textbf{AUROC} (area under the ROC curve, the probability that a random positive scores above a random negative):
\begin{equation}\mathrm{AUROC}=\int_0^1\mathrm{TPR}\;d(\mathrm{FPR})\end{equation}

\textbf{Youden's index and calibrated threshold} ($\tau=-0.0009178519$):
\begin{equation}J(t)=\mathrm{TPR}(t)-\mathrm{FPR}(t),\qquad \tau=-\arg\max_t J(t)\end{equation}

\textbf{Decision rule:}
\begin{equation}\text{label}=\begin{cases}\text{Likely pathogenic}, & \Delta<\tau,\\ \text{Likely benign}, & \Delta\ge\tau\end{cases}\end{equation}

\textbf{Confidence} ($\sigma_{LOF}=0.0015140239$ for a pathogenic label, $\sigma_{FUNC}=0.0009016589$ for a benign label):
\begin{equation}\mathrm{conf}=\min\!\left(1,\;\frac{|\Delta-\tau|}{\sigma_{class}}\right)\end{equation}

\textbf{Mutated sequence length} (deletion of $n$ bases, insertion of $m$ bases) and \textbf{GPU forward passes} ($N$ variants, $U$ unique reference windows):
\begin{equation}|w^{var}|=|w^{ref}|-n\ \text{or}\ |w^{ref}|+m,\qquad \text{passes}_{\text{request}}=2,\qquad \text{passes}_{\text{calibration}}=U+N\end{equation}
\endgroup
\clearpage
\label{s:app:end}
